% ===================================================================
%  तालिका सं. 5 — घर और उसका बनना।
%  Traduction 2026 d'apres texto/io, controlee sur texto/fr et
%  texto/en. Consigne : voir texto/hi/10-talika-01.tex.
%
%  Le tableau 5 est un DIALOGUE, et les deux volumes composent le nom
%  du parleur en petites capitales : on garde \textsc{}, que le rendu
%  laisse tomber en gardant le texte.
%
%  Les renvois du plan sont des LETTRES — (a) le couloir, (b) le
%  salon... — et non des numeros. On les ecrit comme l'ido les ecrit,
%  y compris la ou l'imprimeur a compose « (1) » pour « (l) ».
% ===================================================================

%%K t05-apar-1 apar
\VUsaut{6.0mm}
\VUtitre{15.0pt}{60}{तालिका सं. 5}
\VUsaut{1.4mm}
\VUtitre{11.4pt}{0}{घर और उसका बनना।}
\VUsaut{2.2mm}

%%K t05-tit-1 sub
\VUpk{10.2pt}{40}{अच्छी भेंट।}
\VUsaut{3.0mm}

%%K t05-01-1 p
\textsc{जॉन}। --- चाचा, मैं बहुत जानना चाहता हूँ कि \VUgras{घर} कैसे बनता है।

%%K t05-01-2 p
\textsc{चाचा} \textsuperscript{(80)}। --- यह सहज है, बेटा; मेरे साथ आओ और मैं
तुम्हें वह घर दिखाऊँगा जो श्री बर्नार्ड \textsuperscript{(79)} बनवा रहे हैं और
जिसमें मैं रहने जा रहा हूँ।

%%K t05-01-3 p
\textsc{जॉन}। --- और इस घर का \VUgras{नक़्शा} \textsuperscript{(76)} किसने बनाया?

%%K t05-01-4 p
\textsc{चाचा}। --- \VUgras{वास्तुकार} \textsuperscript{(75)} ने; उन्होंने बहुत
साफ़ चित्र बनाया, वह मान लिया गया, और \VUgras{ठेकेदार} \textsuperscript{(77)} ने
काम आरंभ कर दिया।

%%K t05-01-5 p
\textsc{जॉन}। --- और ठेकेदार कौन हैं?

%%K t05-01-6 p
\textsc{चाचा}। --- श्री श्वेत, तुम्हारे साथी जेम्स के पिता। वे वास्तुकार श्री रू
के साथ मिलकर मज़दूरों को निर्देश देते हैं।

%%K t05-01-7 p
लो! श्री श्वेत ठीक समय पर अपना \VUgras{पैमाना} \textsuperscript{(78)} लिए आ रहे
हैं, और श्री रू अपना नक़्शा लिए।

%%K t05-01-8 p
नमस्ते, महोदय। आप कैसे हैं?

%%K t05-01-9 p
\textsc{श्री रू}। --- बहुत अच्छा, धन्यवाद। नमस्ते जॉन; लड़का कितना बढ़ गया है!

%%K t05-01-10 p
\textsc{चाचा}। --- सचमुच, यह पूरा छोटा आदमी है। यह बहुत गंभीर है; जो कुछ देखता
है, उसे समझाना पड़ता है। आज यह जानना चाहता है कि घर कैसे बनता है।

%%K t05-tit-2 sub
\VUpk{10.2pt}{40}{मज़दूर।}
\VUsaut{3.0mm}

%%K t05-01-11 p
\textsc{श्री श्वेत}। --- तो मेरे साथ चलो, मेरे नन्हे मित्र, और मैं तुम्हें अभी
समझा देता हूँ, क्योंकि मुझे सीखने के इच्छुक बच्चे बहुत भाते हैं।

%%K t05-01-12 p
वह मज़दूर देखते हो, जिसके हाथ में \VUgras{बेलचा} \textsuperscript{(82)} है : वह
\VUgras{खुदाई करने वाला} \textsuperscript{(81)} है। वह पानी की नली बिछाने के लिए
एक \VUgras{खाई} \textsuperscript{(46)} खोद रहा है। उसी ने वह ज़मीन भी खोदी थी
जहाँ घर खड़ा है, ताकि राज चारों \VUgras{दीवारें} उठा सके। उन दीवारों का जो भाग
ज़मीन में है उसे \VUgras{नींव} \textsuperscript{(2)} कहते हैं। नींव से घिरी जगह
\VUgras{तहख़ाना} \textsuperscript{(3)} बनाती है। इस तहख़ाने में एक छोटी खिड़की है
जिसे \VUgras{हवादान} \textsuperscript{(5)} कहते हैं, एक \VUgras{दरवाज़ा}
\textsuperscript{(4)} और एक सीढ़ी।

%%K t05-01-13 p
\VUgras{राज} \textsuperscript{(108)} दीवारें उठाता गया, \VUgras{दरवाज़ों}
\textsuperscript{(8)} और \VUgras{खिड़कियों} \textsuperscript{(17)} के लिए जगह
छोड़ता हुआ।

%%K t05-01-14 p
\textsc{जॉन}। --- पर मुझे तो यह राज दिखाई ही नहीं देता!

%%K t05-01-15 p
\textsc{श्री श्वेत}। --- उसका घर का काम अब पूरा हो चुका है। वह उधर
\VUgras{अस्तबल} \textsuperscript{(54)} के पास है, अपने \VUgras{मचान}
\textsuperscript{(110)} पर चढ़ा, \VUgras{धोबीघर} \textsuperscript{(56)} की दीवार
बना रहा है।

%%K t05-01-16 p
उसके पास एक करनी, एक नाँद और एक \VUgras{साहुल} \textsuperscript{(109)} है। पर ये
नाम तुम्हें याद न रहेंगे।

%%K t05-01-17 p
\textsc{जॉन}। --- क्यों नहीं रहेंगे, मैं चाचा से एक छोटी करनी और एक नाँद
माँगूँगा, और साहुल तो मैं सुतली के टुकड़े से आप ही बना लूँगा।

%%K t05-tit-3 sub
\VUsaut{3.0mm}
\VUpk{10.2pt}{40}{सामग्री।}
\VUsaut{3.0mm}

%%K t05-01-18 p
\textsc{चाचा}। --- वाह, बहुत अच्छे। पर बताओ तो जॉन, घर बनाने के लिए क्या चाहिए?

%%K t05-01-19 p
\textsc{जॉन}। --- \VUgras{तराशा हुआ पत्थर} \textsuperscript{(59)} और
\VUgras{गारा} \textsuperscript{(65)} चाहिए।

%%K t05-01-20 p
\textsc{चाचा}। --- ठीक। उस बड़े हैट वाले आदमी को देखो, अपनी \VUgras{गाड़ी}
\textsuperscript{(136)} पर। वह \VUgras{गाड़ीवान} \textsuperscript{(135)} है। वह हर
दिन यहाँ \VUgras{पत्थर के ढोके} \textsuperscript{(60)} लाता है। वह अपनी गाड़ी
आँगन में ख़ाली करता है, और \VUgras{संगतराश} \textsuperscript{(83)} ढोकों को
तराशते और अपनी \VUgras{पत्थर-आरी} \textsuperscript{(85)} से चीरते हैं। एक
\VUgras{मज़दूर} \textsuperscript{(146)} उन्हें राज तक पहुँचाता है, जो उन्हें अपनी
जगह बिठाता है।

%%K t05-01-21 p
इस बीच \VUgras{गारा बनाने वाला} \textsuperscript{(87)} गारा तैयार करता है। उसे
\VUgras{बालू} \textsuperscript{(62)}, \VUgras{चूना} \textsuperscript{(63)} और
\VUgras{पानी} \textsuperscript{(64)} चाहिए। चूना उसके पास एक \VUgras{बोरी}
\textsuperscript{(71)} में है; एक \VUgras{राज-सहायक} \textsuperscript{(111)} उसे
\VUgras{ठेले} \textsuperscript{(112)} में बालू लाता है, और \VUgras{शागिर्द}
\textsuperscript{(113)} पंप \textsuperscript{(58)} पर है। वह एक \VUgras{बाल्टी}
\textsuperscript{(114)} भर रहा है। गारा जल्दी ही तैयार हो जाएगा।
\VUgras{गारा ढोने वाला} \textsuperscript{(107)} उसे लेकर सीढ़ी से मचान पर चढ़ा ले
जाता है, जहाँ एक राज घर का वह पाखा पूरा कर रहा है।

%%K t05-01-22 p
और अब बताओ, जो मज़दूर \VUgras{छत} \textsuperscript{(31)} बनाता है उसे क्या कहते
हैं?

%%K t05-01-23 p
\textsc{जॉन}। --- वह \VUgras{छत बनाने वाला} \textsuperscript{(118)} है।

%%K t05-tit-4 sub
\VUsaut{3.0mm}
\VUpk{10.2pt}{40}{घर की छत।}
\VUsaut{3.0mm}

%%K t05-01-24 p
\textsc{श्री श्वेत}। --- हाँ, पर पहले \VUgras{बढ़ई} \textsuperscript{(115)} आता
है।

%%K t05-01-25 p
बढ़ई \VUgras{छत का ढाँचा} \textsuperscript{(35)} बनाता है। वह \VUgras{कड़ियों}
\textsuperscript{(37)} को \VUgras{खूँटों} \textsuperscript{(117)} से जोड़ता है। ऊपर
वे मज़दूर, अपनी चौड़ी मख़मली निकर में, बढ़ई हैं। एक छोटी कड़ी थामे है और दूसरा
अपने \VUgras{कुल्हाड़े} \textsuperscript{(116)} से एक खूँटा गढ़ रहा है।
\VUgras{चिमनी} \textsuperscript{(41)} के पास वाला छत बनाने वाला है। वह
\VUgras{शहतीरों} (*) पर पट्टियाँ ठोंकता है, और पट्टियों पर \VUgras{स्लेटें}
\textsuperscript{(66)}। कभी-कभी वह खपरैल बिछाता है।

%%K t05-noto-1 noto
\VUnotes{143.2mm}{%
(*) \VUgras{Balk-o} = जर्मन \textit{Balken}, अंग्रेज़ी \textit{joist},
फ़्रेंच \textit{solive}, इतालवी \textit{travicello}, रूसी \textit{balka},
स्पेनी और पुर्तगाली \textit{viga}। एक तकनीकी धातु जो अभी अपनाई नहीं गई, पर यहाँ
फ़्रेंच शब्द \textit{solive} का अर्थ देने के लिए सुझाई जाती है, जो न
\textit{trabo} (फ़्रेंच \textit{poutre}) है न छोटी कड़ी। वह चौकोर गढ़ी लकड़ी है
जो एक फ़र्श और एक छत सँभालती है।%
}

%%K t05-01-26 p
अस्तबल की छत \VUgras{खपरैल की} \textsuperscript{(55)} है, घर की
\VUgras{स्लेट की} \textsuperscript{(32)}, और बरामदे की \VUgras{जस्ते की}
\textsuperscript{(24)}।

%%K t05-01-27 p
\textsc{जॉन}। --- क्या जस्ते की छतें भी छत बनाने वाला ही बनाता है?

%%K t05-01-28 p
\textsc{श्री श्वेत}। --- नहीं, वह \VUgras{जस्तासाज़} \textsuperscript{(121)} या
\VUgras{नलसाज़} का काम है — वही जो घर की छत में एक \VUgras{रोशनदान-खिड़की}
\textsuperscript{(38)} बिठा रहा है। वही \VUgras{छत के परनाले} \textsuperscript{(43)}
और \VUgras{उतरती नलियाँ} \textsuperscript{(42)} भी लगाता है। वह जस्ते और टीन को
टाँका लगाता है। उसी ने \VUgras{तड़ित-चालक} \textsuperscript{(40)} और \VUgras{वात-सूचक}
\textsuperscript{(39)} \VUgras{मुँडेर} \textsuperscript{(36)} पर जड़े थे।

%%K t05-01-29 p
\textsc{जॉन}। --- तो क्या घर पूरा हो गया?

%%K t05-tit-5 sub
\VUsaut{3.0mm}
\VUpk{10.2pt}{40}{औज़ार।}
\VUsaut{3.0mm}

%%K t05-01-30 p
\textsc{श्री श्वेत}। --- पूरा नहीं। \VUgras{पलस्तर करने वाला}
\textsuperscript{(99)} \VUgras{ईंटों} \textsuperscript{(61)} से वे दीवारें बनाता है
जो घर को अलग-अलग कमरों में बाँटती हैं। वह \VUgras{पलस्तर} \textsuperscript{(70)} \VUgras{छतों}
\textsuperscript{(26)} और दीवारों पर चढ़ाता है। इस समय वह
\VUgras{बरामदे} \textsuperscript{(33)} की छत पर काम कर रहा है। राज की तरह उसके पास
भी एक \VUgras{करनी} \textsuperscript{(100)} और एक \VUgras{नाँद}
\textsuperscript{(101)} है।

%%K t05-01-31 p
फिर बढ़ई दरवाज़े, खिड़कियाँ, \VUgras{लकड़ी के फ़र्श} \textsuperscript{(16)} और
\VUgras{झिलमिलियाँ} \textsuperscript{(19)} बनाता है।

%%K t05-01-32 p
इस समय वह आँगन में अपनी \VUgras{कारीगरी की मेज़} \textsuperscript{(90)} पर काम कर
रहा है। उसके पास एक \VUgras{आरी} \textsuperscript{(94)} है, \VUgras{लकड़ी}
\textsuperscript{(69)} चीरने के लिए, एक \VUgras{रंदा} \textsuperscript{(91)}, एक
\VUgras{शिकंजा} \textsuperscript{(92)}, एक \VUgras{छेनी}
\textsuperscript{(94 \textit{bis})}, एक \VUgras{परकार} \textsuperscript{(95)} और
लकड़ी का एक \VUgras{मुँगरा} \textsuperscript{(95 \textit{bis})}।

%%K t05-01-33 p
\textsc{जॉन}। --- अरे, तब तो बढ़ई होना राज होने से भी अधिक मज़े का है। चाचा, आप
मुझे एक कारीगरी की मेज़, एक आरी और एक रंदा दिलाएँगे? मैं मेदोर के लिए एक घर
बनाऊँगा।

%%K t05-01-34 p
\textsc{चाचा}। --- हाँ बेटा।

%%K t05-01-35 p
\textsc{जॉन}। --- कितनी ख़ुशी हुई! और वह आदमी कौन है जो सामने के दरवाज़े के पास
खड़ा है?

%%K t05-tit-6 sub
\VUsaut{3.0mm}
\VUpk{10.2pt}{40}{रँगाई।}
\VUsaut{3.0mm}

%%K t05-01-36 p
\textsc{श्री श्वेत}। --- वह \VUgras{रँगसाज़} \textsuperscript{(102)} है। वह अपनी
सीढ़ी पर खड़ा है, \VUgras{रंग} \textsuperscript{(103)} की एक हाँडी और अपना
\VUgras{ब्रश} \textsuperscript{(104)} लिए। वह सामने के दरवाज़े की लकड़ी पर रंग
चढ़ा रहा है ताकि वह अधिक टिके।

%%K t05-01-37 p
कमरों में दीवार का काग़ज़ भी वही चिपकाता है।

%%K t05-01-38 p
\VUgras{सजावट करने वाला} \textsuperscript{(107)}, जो \VUgras{सीढ़ी}
\textsuperscript{(106)} पर चढ़ा \VUgras{बालाख़ाने} \textsuperscript{(28)} पर
है, एक \VUgras{परदा} \textsuperscript{(29)} लगा रहा है।

%%K t05-01-39 p
बड़ी खिड़की के पास खड़ा मज़दूर \VUgras{शीशासाज़} \textsuperscript{(96)} है। वह
\VUgras{शीशे} \textsuperscript{(18)} \VUgras{मसाले} \textsuperscript{(73)} से
जमाता है। वह दूसरा मज़दूर, जो तहख़ाने के दरवाज़े के पास घुटनों के बल है,
\VUgras{तालासाज़} \textsuperscript{(97)} है। वह एक \VUgras{ताला}
\textsuperscript{(6)} बिठा रहा है। उसकी \VUgras{रेती} \textsuperscript{(98)} उसके
पास पड़ी है।

%%K t05-01-40 p
\textsc{जॉन}। --- जो आदमी अपने \VUgras{हथौड़े} \textsuperscript{(84)} से लोहे के
टुकड़े पर चोट कर रहा है, क्या वह भी तालासाज़ है?

%%K t05-01-41 p
\textsc{श्री श्वेत}। --- ठीक कहें तो वह \VUgras{लोहार} \textsuperscript{(125)} है।
वह \VUgras{लोहे का सामान} \textsuperscript{(67)} गढ़ रहा है
\VUgras{बिजलीवाले} \textsuperscript{(124)} के लिए, जो \VUgras{गाड़ीख़ाने}
\textsuperscript{(51)} की छत पर है। उसके पास एक \VUgras{भट्ठी} \textsuperscript{(126)}, एक \VUgras{निहाई}
\textsuperscript{(127)} और एक \VUgras{संडासी} \textsuperscript{(128)} है।

%%K t05-01-42 p
बिजलीवाला टेलीफ़ोन का \VUgras{ताँबे का तार} \textsuperscript{(44)} जोड़ रहा है।

%%K t05-tit-7 sub
\VUpk{10.2pt}{40}{बाग़वानी।}
\VUsaut{3.0mm}

%%K t05-01-43 p
\textsc{चाचा}। --- \VUgras{आँगन} \textsuperscript{(45)} के सिरे पर,
\VUgras{जँगले} \textsuperscript{(47)} और \VUgras{गाड़ी-फाटक}
\textsuperscript{(48)} के पास, तुम्हें \VUgras{माली} \textsuperscript{(129)} दिखाई
देता है। वह छोटा \VUgras{बग़ीचा} \textsuperscript{(49)} तैयार कर रहा है। उसने
अपने \VUgras{फावड़े} \textsuperscript{(131)} से मिट्टी खोद डाली है, वह बीज बो रहा
है और \VUgras{पंजे} \textsuperscript{(130)} से समतल कर रहा है। फिर अपने
\VUgras{झारे} \textsuperscript{(132)} से वह \VUgras{क्यारियों}
\textsuperscript{(150)} को सींचेगा और पौधे उगेंगे।

%%K t05-01-44 p
\VUgras{साईस} \textsuperscript{(133)}, जो अभी घोड़े की देखभाल कर उसे खिला चुका
है, \VUgras{बग्घी} \textsuperscript{(52)} साफ़ कर रहा है — स्पंज, एक
\VUgras{हाथ-पंप} \textsuperscript{(134)} और पानी की बाल्टी से। फिर वह उसे गाड़ीख़ाने में
रख देगा और भारी \VUgras{अर्गला} \textsuperscript{(53)} से दरवाज़ा बंद कर देगा।

%%K t05-01-45 p
लो, अब तो संतोष हुआ होगा; समझाना बहुत हो चुका।

%%K t05-01-46 p
\textsc{जॉन}। --- हाँ चाचा, पर मैं घर के भीतर भी देखना चाहूँगा।

%%K t05-01-47 p
\textsc{चाचा}। --- वह कल के लिए। श्री रू तुम्हें नक़्शा समझाएँगे और हर कमरे का
नाम बताएँगे।

%%K t05-tit-8 sub
\VUsaut{3.0mm}
\VUpk{10.2pt}{40}{घर की भीतरी बनावट।}
\VUsaut{2.4mm}

%%K t05-apar-2 apar
{\centering\textit{(नक़्शा देखिए।)}\par}
\VUsaut{2.4mm}

%%K t05-01-48 p
\textsc{वास्तुकार}। --- आओ जॉन, मैं तुम्हें घर के अलग-अलग भाग दिखाता हूँ।

%%K t05-01-49 p
चलो \VUgras{भूतल} \textsuperscript{(7)} में चलें। यह रहा \VUgras{घंटी} का बटन
\textsuperscript{(11)}। हम तीन \VUgras{सीढ़ियाँ} \textsuperscript{(14)} चढ़ते हैं, जो
\VUgras{ड्योढ़ी} \textsuperscript{(9)} की हैं, फिर \VUgras{देहरी}
\textsuperscript{(10)} पार करते हैं, जो \VUgras{सामने के दरवाज़े}
\textsuperscript{(8)} की है, और यह
रहे हम \VUgras{ज़ीने} \textsuperscript{(13)} के नीचे।

%%K t05-01-50 p
\textsc{जॉन}। --- यह गलियारा है।

%%K t05-01-51 p
\textsc{वास्तुकार}। --- नहीं, यह \VUgras{ड्योढ़ी-कक्ष} \textsuperscript{(12)} है।
\VUgras{गलियारा} \textsuperscript{(a)} सिरे पर है। दाहिनी ओर यह रहे
\VUgras{बैठक} \textsuperscript{(b)} और \VUgras{भोजन-कक्ष} \textsuperscript{(c)};
भोजन-कक्ष के सामने \VUgras{रसोई} \textsuperscript{(d)} और \VUgras{भंडार-कक्ष}
\textsuperscript{(e)}; फिर आगे की ओर \VUgras{अध्ययन-कक्ष} \textsuperscript{(f)}।
\VUgras{बरामदा} \textsuperscript{(23)}, अपने \VUgras{काँच के दरवाज़े}
\textsuperscript{(25)} समेत, बैठक के पास है।

%%K t05-01-52 p
अब हम ज़ीना चढ़ेंगे। \VUgras{जंगला} \textsuperscript{(15)} थामे रहो और सीढ़ियाँ
गिनो। चौदह हैं। यह रही \VUgras{चौकी} \textsuperscript{(m)} : हम पहली
\VUgras{मंज़िल} \textsuperscript{(27)} पर हैं।

%%K t05-tit-9 sub
\VUsaut{3.0mm}
\VUpk{10.2pt}{40}{पहली मंज़िल।}
\VUsaut{3.0mm}

%%K t05-01-53 p
ज़ीने के सामने \VUgras{शौचालय} \textsuperscript{(20)} और \VUgras{वेश-कक्ष}
\textsuperscript{(j)} हैं। दाहिनी ओर एक सुंदर \VUgras{शयन-कक्ष}
\textsuperscript{(g)}, जिसकी दो खिड़कियाँ घर के \VUgras{सामने} \textsuperscript{(1)}
की ओर खुलती हैं। बाईं ओर \VUgras{स्नान-कक्ष} \textsuperscript{(1)} और
\VUgras{अतिथि-कक्ष} \textsuperscript{(i)} हैं, जिसमें एक बड़ा \VUgras{आला}
\textsuperscript{(h)} है। शयन-कक्ष के पास \VUgras{बच्चों का कमरा}
\textsuperscript{(k)} है। वह एक चौड़ी \VUgras{छत-बारजा} \textsuperscript{(21)} में
खुलता है। इस बारजे के नीचे एक और शौचालय \textsuperscript{(20)} है, और दीवार से
सटी यह रही \VUgras{घंटी} \textsuperscript{(22)}, जो भोजन के लिए बजाई जाती है।

%%K t05-01-54 p
\textsc{जॉन}। --- चलिए \VUgras{दूसरी मंज़िल} पर चलें।

%%K t05-01-55 p
\textsc{वास्तुकार}। --- दूसरी मंज़िल है ही नहीं। छत के नीचे केवल
\VUgras{अटारियाँ} \textsuperscript{(33)} और भंडार \textsuperscript{(34)} हैं। हमारा
चक्कर पूरा हुआ, अब नीचे उतर सकते हैं।

%%K t05-01-56 p
\textsc{जॉन}। --- धन्यवाद महोदय, बहुत रोचक रहा। मैं अपने पिता को सब कुछ बताऊँगा
जो मैंने देखा।
\VUsaut{4.0mm}
\VUfilet{22mm}
